% ============================================================
% PRN232 - Section 7: Testing (Vietnamese)
% ============================================================
\section{\bi{Kiểm thử và đánh giá}{Testing and Evaluation}}
\label{sec:test}

\subsection{\bi{Chiến lược kiểm thử}{Test Strategy}}

\bi{
Dự án áp dụng chiến lược kiểm thử theo tầng, trong đó trọng tâm là
kiểm thử đơn vị cho logic nghiệp vụ vì đây là phần chứa nhiều quy tắc
tinh vi nhất. Điểm đặc biệt của dự án là \emph{toàn bộ bộ kiểm thử
back-end chạy được mà không cần cơ sở dữ liệu}: mỗi spec file tiêm
một đối tượng Prisma giả (fake client), nhờ vậy chu kỳ kiểm thử rất
ngắn và hoàn toàn xác định.
}{
The project applies a layered test strategy, with unit testing of the
business rules taking priority because that layer holds the most
subtle logic. A distinctive property of this project is that the
\emph{entire back-end suite runs without a database}: each spec file
injects a fake Prisma client, so the test cycle is short and fully
deterministic.
}

\subsection{\bi{Kết quả thực thi}{Execution Results}}

\bi{
Các bộ kiểm thử được chạy thật trên máy dự án ngày 28 tháng 10 năm
2026. Toàn bộ 482 kiểm thử đều đạt, không có kiểm thử nào bị bỏ qua
hoặc lỗi.
}{
All suites were executed on the project machine on 28 October 2026.
All 482 tests pass, with no skipped or errored test.
}

\begin{table}[t]
\caption{\bi{Kết quả kiểm thử tự động của ba bộ kiểm thử.}{Automated
test results across the three suites.}}
\label{tab:tests}
\centering
\footnotesize
\fitcol{% Auto-generated from deliverables/02_Testing/test-results - do not edit by hand.
\begin{tabular}{@{}llrr@{}}
\toprule
\textbf{\bi{Bộ kiểm thử}{Test suite}} & \textbf{\bi{Lệnh chạy}{Command}} & \textbf{\bi{Tổng}{Total}} & \textbf{\bi{Đạt}{Passed}} \\
\midrule
\bi{Backend (node:test)}{Backend (node:test)} & \texttt{npm test} & 163 & 163 \\
\bi{Frontend (Vitest)}{Frontend (Vitest)} & \texttt{npx vitest run} & 228 & 228 \\
\bi{Mobile (flutter test)}{Mobile (flutter test)} & \texttt{flutter test} & 91 & 91 \\
\midrule
\textbf{\bi{Tổng cộng}{Total}} & & \textbf{482} & \textbf{482} \\
\bottomrule
\end{tabular}
}
\end{table}

\subsection{\bi{Bộ kiểm thử back-end}{Back-End Test Suite}}

\bi{
Bộ kiểm thử back-end dùng trình chạy \texttt{node:test} tích hợp sẵn
của Node.js 20, đi kèm một runner tự viết
(\texttt{backend/test/runner.js}) biên dịch các spec bằng
\texttt{tsc} rồi thực thi tuần tự. Việc chọn trình chạy tích hợp
thay vì Jest hay Vitest giúp loại bỏ hoàn toàn phụ thuộc bên thứ ba
trong bộ kiểm thử, giảm bề mặt lỗi và giữ cho quy trình CI ổn định.
}{
The back-end suite uses Node.js 20's built-in \texttt{node:test}
runner together with a custom runner (\texttt{backend/test/runner.js})
that compiles the specs with \texttt{tsc} and executes them
sequentially. Choosing the built-in runner over Jest or Vitest removes
every third-party dependency from the test path, shrinking the failure
surface and keeping the CI pipeline stable.
}

\begin{table}[t]
\caption{\bi{Phân bố kiểm thử back-end theo tệp spec.}{Back-end test
distribution per spec file.}}
\label{tab:backend-specs}
\centering
\scriptsize
\begin{tabular}{@{}lrl@{}}
\toprule
\textbf{Tệp spec} & \textbf{\bi{Kiểm thử đạt}{Tests passed}} & \textbf{\bi{Tỉ lệ đạt}{Pass rate}} \\
\midrule
\texttt{ai-curator.service.spec.ts}       & 19 & 100\% \\
\texttt{epg-filler.spec.ts}               & 18 & 100\% \\
\texttt{live-studio.service.spec.ts}      & 18 & 100\% \\
\texttt{auth.service.spec.ts}             & 17 & 100\% \\
\texttt{ai-curator.e2e.spec.ts}           & 13 & 100\% \\
\texttt{watchlist.service.spec.ts}        & 11 & 100\% \\
\texttt{ingest-sources.spec.ts}           & 11 & 100\% \\
\texttt{ai-curator-stats.service.spec.ts} & 10 & 100\% \\
\texttt{channels.service.spec.ts}         &  7 & 100\% \\
\texttt{programs.conflict.spec.ts}        &  7 & 100\% \\
\texttt{programs.counter.spec.ts}         &  8 & 100\% \\
\texttt{audit-logger.service.spec.ts}     &  8 & 100\% \\
\texttt{epg-timeline.spec.ts}             &  5 & 100\% \\
\texttt{search.service.spec.ts}           &  4 & 100\% \\
\texttt{users.service.spec.ts}            &  4 & 100\% \\
\texttt{billing.service.spec.ts}          &  3 & 100\% \\
\midrule
\textbf{Tổng} & \textbf{163} & \textbf{100\%} \\
\bottomrule
\end{tabular}
\end{table}

\bi{
Hai tệp kiểm thử đáng chú ý nhất vì chúng bảo vệ trực tiếp các thuật
toán nghiệp vụ cốt lõi. \texttt{programs.conflict.spec.ts} kiểm tra
thuật toán phát hiện xung đột trên cả ba dạng giao thoa, trường hợp
giá trị thời lượng bằng \texttt{NULL}, và trường hợp loại trừ chính
sự kiện đang cập nhật. \texttt{epg-filler.spec.ts} kiểm tra tính
đầy đủ 24 giờ của lưới, tính xác định của lựa chọn ngẫu nhiên có
hạt giống, giới hạn số lần phát lại, và hành vi tách tập.
}{
Two spec files deserve particular mention because they directly
protect the core algorithms. \texttt{programs.conflict.spec.ts}
exercises the conflict detector across all three overlap forms, the
\texttt{NULL}-duration case and the self-exclusion case used when
updating an event. \texttt{epg-filler.spec.ts} verifies the
completeness of the 24-hour grid, the determinism of the seeded
selection, the repeat limit and the episode-splitting behaviour.
}

\subsection{\bi{Bộ kiểm thử web}{Web Test Suite}}

\bi{
228 kiểm thử Vitest trong 76 bộ kiểm thử bao phủ cả ba loại đối
tượng: hàm thuần, hook React và thành phần giao diện. Các hàm thuần
được kiểm thử với đầy đủ trường hợp biên (chuỗi rỗng, không phân biệt
hoa thường, tham số không hợp lệ); các hook được kiểm thử với trạng
thái tải, thành công và lỗi; các thành phần được kiểm thử bằng
\texttt{React Testing Library} với các truy vấn theo vai trò người
dùng thay vì kiểm tra cấu trúc DOM thô.
}{
228 Vitest tests in 76 suites cover three kinds of subject: pure
functions, React hooks and interface components. Pure functions are
tested at their boundaries (empty strings, case-insensitivity,
invalid parameters); hooks are tested across loading, success and
error states; components are tested with \texttt{React Testing
Library} using user-facing queries rather than raw DOM structure
assertions.
}

\subsection{\bi{Bộ kiểm thử di động và kiểm thử E2E}{Mobile Suite and End-to-End Tests}}

\bi{
Bộ kiểm thử Flutter gồm 91 kiểm thử trong 8 bộ, tổ chức theo cùng
cấu trúc thư mục với mã nguồn: \texttt{test/data},
\texttt{test/logic}, \texttt{test/presentation}. Bộ này chạy với
\texttt{flutter test --coverage}, sinh báo cáo phủ \texttt{lcov.info}
được tải lên dưới dạng artefact của CI.

Bên cạnh đó, một bộ kiểm thử E2E dựa trên Playwright \cite{test}
thực hiện tám kịch bản trên trình duyệt Chromium: biểu mẫu đăng nhập
từ chối email không hợp lệ, mật khẩu sai hiển thị thông báo lỗi,
trang EPG không lỗi, trang kênh hiển thị nội dung, trang tìm kiếm
nhận và hiển thị kết quả, trang chủ hiển thị mặt đất công khai, và hai
kịch bản kiểm tra hành vi chuyển hướng khi chưa đăng nhập.
}{
The Flutter suite contains 91 tests in 8 suites, organised with the
same directory layout as the source: \texttt{test/data},
\texttt{test/logic}, \texttt{test/presentation}. It runs with
\texttt{flutter test --coverage} and produces an \texttt{lcov.info}
report uploaded as a CI artefact.

Alongside it, a Playwright-based end-to-end suite \cite{test}
executes eight scenarios in Chromium: the login form rejects an
invalid email, a wrong password surfaces an error toast, the EPG page
renders without crashing, the channels page displays content, the
search page accepts input and renders results, the root URL shows a
public landing surface, and two scenarios verify unauthenticated
redirect behaviour.
}

\begin{keybox}{\bi{Nhận xét về hai lỗi E2E còn lại}{Note on the two
remaining E2E failures}}
\bi{
Trong lần chạy E2E gần nhất, sáu trong tám kịch bản đạt và hai kịch
bản thất bại: một kịch bản đăng nhập STAFF không thoát khỏi
\texttt{/login} trong thời hạn chờ, và một kịch bản chuyển hướng từ
\texttt{/studio} khi chưa đăng nhập. Cả hai đều phụ thuộc vào dữ liệu
seed sẵn có của môi trường kiểm thử chứ không phải vào lỗi logic của
ứng dụng: kịch bản thứ hai nhận mã trạng thái HTTP 401 thay vì bị
chuyển hướng, tức là chính sách bảo vệ quyền đã hoạt động đúng nhưng
điều kiện kiểm thử chỉ chấp nhận một trong ba dạng phản hồi. Kịch
bản thứ nhất phụ thuộc tài khoản STAFF có sẵn trong cơ sở dữ liệu
kiểm thử. Các trường hợp này được ghi nhận đầy đủ trong
\texttt{TEST\_EXECUTION\_REPORT.md} đi kèm.
}{
In the most recent end-to-end run, six of eight scenarios passed and
two failed: a STAFF login that did not leave \texttt{/login} within
the wait timeout, and an unauthenticated redirect from \texttt{/studio}.
Both depend on seeded data in the test environment rather than on an
application logic defect: the second scenario received an HTTP 401
instead of a redirect, meaning the authorisation policy behaved
correctly while the assertion accepted only one of three response
forms. The first depends on a STAFF account existing in the test
database. Both are documented in full in the accompanying
\texttt{TEST\_EXECUTION\_REPORT.md}.
}
\end{keybox}

\subsection{\bi{Đánh giá chất lượng}{Quality Assessment}}

\bi{
Dựa trên kết quả trên, hệ thống đạt mức độ tin cậy phù hợp với môi
trường triển khai giáo dục. Điểm mạnh rõ nhất là tỉ lệ phủ của logic
nghiệp vụ: hai thuật toán rủi ro nhất (phát hiện xung đột và điền
lịch EPG) đều được kiểm thử bằng các trường hợp biên thay vì chỉ kiểm
thử đường hạnh phúc. Điểm cần cải thiện là mức độ phủ của tầng E2E,
hiện chỉ gồm tám kịch bản trình duyệt và chưa bao gồm kịch bản phát
sóng trực tiếp vốn đòi hỏi phần cứng và cơ sở hạ tầng MediaMTX.
}{
Based on these results, the system reaches a reliability level
appropriate for an educational deployment. The clearest strength is
coverage of the business logic: the two riskiest algorithms (conflict
detection and EPG filling) are tested at their boundaries rather than
only on the happy path. The area needing improvement is end-to-end
coverage, currently limited to eight browser scenarios and excluding
the live-broadcast scenarios that would require dedicated hardware and
MediaMTX infrastructure.
}
